\documentclass[11pt,a4paper]{article}
\usepackage[margin=1in]{geometry}
\usepackage{hyperref}
\usepackage{enumitem}
\usepackage{graphicx}
\usepackage{xcolor}

% -----------------------------------------------------
% bvraghav-tiet-ucs503p-article.sty
% -----------------------------------------------------

% Variable: Lab Instructor
\makeatletter \newcommand{\BvrSetLabInstructor}[1]{%
  \gdef\Bvr@LabInstructor{#1}}
% default
\newcommand{\Bvr@LabInstructor}{Dr. Raghav %
  B. Venkataramaiyer}
% public accessor
\newcommand{\BvrLabInstructorValue}{\Bvr@LabInstructor}
\makeatother

% Add subtitle
\usepackage{titling}
% -- Set title bold
\pretitle{\begin{center}\bfseries\fontsize{18bp}{18bp}\selectfont}
\makeatletter
\newcommand{\BvrSubtitle}[1]{%
  \posttitle{%
    \par\end{center}
    \begin{center}\large#1\end{center}
    \vskip 0.5em}%
}
\makeatother

% Hints in the section title(s)
\newcommand{\BvrHint}[1]{%
  {\normalfont\normalsize\itshape\hfill #1}}

% -----------------------------------------------------

\title{ChargEV: Find. Charge. Go Green.}

\BvrSetLabInstructor{Nisha Thakur}

\BvrSubtitle{%
  The Smart EV Charging Station Platform \\
  Project Proposal for UCS503P  \\
  Submitted to: \BvrLabInstructorValue \\ \bfseries
  Thapar Institute of Engineering and Technology% 
}

\author{
  Pehal, CSED%
  \thanks{Roll No: \texttt{1024170103}, %
    Email: \texttt{ppehal\_be24\@thapar.edu}}
  \and 
  Briti Arora, CSED%
  \thanks{Roll No: \texttt{1024170431}, %
    Email: \texttt{barora2\_be24\@thapar.edu}}
}

\date{15 August 2026}

\begin{document}
\maketitle

\section{Higher-order goal%
  \BvrHint{\ldots secondary framing}}
This project aims to accelerate the adoption of electric
vehicles (EVs) by removing a major piece of charging
infrastructure friction: the uncertainty around where to
charge, whether a station is available, and whether a slot
can be guaranteed. While the broader goal is a smarter,
greener transportation ecosystem, the immediate engineering
objective is to translate \emph{predictable, low-anxiety EV
charging} into a measurable, deployable software solution.

\section{Time-to-value %
\BvrHint{\ldots primary heuristic}}
\begin{itemize}[leftmargin=0.7cm]
    \item We will deliver meaningful increments quickly by prototyping the core discover-book-charge workflow and validating it with users within a short iteration window.
    \item We will prioritize fast feedback over perfection by using weekly iterations (and targeting CI-backed automation early) to reduce release friction.
    \item Success is defined by what can be delivered and measured in the first iteration -- a working station locator and booking flow -- then improved based on user feedback.
\end{itemize}

\section{Problem Statement}
EV drivers routinely face friction that turns charging into
a source of anxiety rather than a routine task. Common pain
points include:
\begin{itemize}[leftmargin=0.7cm]
    \item \textbf{Discovery gap:} it is hard to find nearby charging stations from a single, reliable source.
    \item \textbf{Unclear availability:} station status (free, occupied, out of service) is rarely known before arriving.
    \item \textbf{Wasted time:} drivers wait at busy stations with no way to reserve a slot in advance.
    \item \textbf{Generic search results:} existing tools do not filter by connector type, distance, availability, or price together.
    \item \textbf{Difficult trip planning:} long-distance EV travel requires manually chaining together charging stops.
    \item \textbf{Low engagement:} there is little incentive for drivers to keep returning to a single platform.
\end{itemize}

\textbf{Why this matters now (contextual relevance):} as EV
ownership grows, even small improvements in charging
predictability translate into measurable reductions in
range anxiety and wasted driver time.

\section{Proposed Solution %
\BvrHint{\ldots Software Proposition}}
\subsection{Overview}
Build a \textbf{web-based} EV charging platform,
\emph{ChargEV}, with the following components:
\begin{itemize}[leftmargin=0.7cm]
    \item \textbf{Centralized station locator:} a single source of truth for nearby EV charging stations.
    \item \textbf{Real-time availability dashboard:} live status per station/connector.
    \item \textbf{Online slot booking \& reservation system:} drivers reserve a charging slot in advance.
    \item \textbf{Smart search \& filters:} filter by connector type, distance, availability, and price.
    \item \textbf{Trip planning with charging stops:} route-aware suggestions for long-distance travel.
    \item \textbf{Rewards, badges \& eco-impact tracking:} gamified engagement layer (``EcoCoins'').
\end{itemize}

\subsection{Core workflow}
The user journey is organized into three phases:
\begin{enumerate}[leftmargin=0.7cm]
    \item \textbf{Discover:} access the cross-device web platform, search/filter nearby stations, and evaluate details such as live availability, pricing, connector compatibility, and reviews.
    \item \textbf{Act:} book a time slot and confirm instantly, then get seamless routing/navigation to the station for a guaranteed charge.
    \item \textbf{Engage:} accumulate EcoCoins and track personal savings and eco-impact, incentivizing return usage.
\end{enumerate}

\subsection{Operational constraints for an educational setting}
\begin{itemize}[leftmargin=0.7cm]
    \item Keep the solution usable on low-to-mid range devices via responsive web design.
    \item Avoid dependence on advanced research algorithms; focus on workflow, real-time data, and system correctness.
    \item Design for deployability using common web hosting, managed databases, and existing location-intelligence APIs.
\end{itemize}

\section{Solution Approach %
\BvrHint{\ldots Engineering focus}}
This is an engineering course project, so we focus on a
practical, scalable web architecture rather than novel
algorithm research.

\subsection{System architecture %
\BvrHint{\ldots web-based solution}}
A typical 3-tier web architecture:
\begin{itemize}[leftmargin=0.7cm]
    \item \textbf{Frontend (the interface):} built for speed and cross-device responsiveness using HTML5, CSS3, Bootstrap~5, and JavaScript.
    \item \textbf{Backend (the engine):} high-performance, event-driven request routing using Node.js and Express.js.
    \item \textbf{Data \& logistics (the brains):} cloud-native storage via MongoDB Atlas, combined with location-intelligence APIs -- Google Maps API, Google Places API, and Google Directions API -- for search, routing, and trip planning.
\end{itemize}

\subsection{CI/CD and fast delivery enabler}
CI/CD will be used to:
\begin{itemize}[leftmargin=0.7cm]
    \item Automatically build and run tests on every change.
    \item Produce repeatable deployment artifacts to staging.
    \item Enable safe and frequent releases of incremental features (e.g., filters, then trip planning, then rewards).
\end{itemize}

\section{Evaluation Criterion %
\BvrHint{\ldots measurable and attributable}}
We define success using measurable metrics tied to the core
discover-book-charge workflow.

\subsection{Primary evaluation metric}
\textbf{Time-to-charge (TTC):} time between a user opening
the app to find a station and successfully starting a
charge at a booked slot.
\begin{itemize}[leftmargin=0.7cm]
    \item Target: median TTC reduced compared to a no-booking baseline (walk-in search).
    \item Attribution: measured from application timestamps (search start vs. booking confirmation vs. charge start).
\end{itemize}

\subsection{Secondary evaluation metrics}
\begin{itemize}[leftmargin=0.7cm]
    \item \textbf{Booking conversion rate:} percentage of searches that convert into a confirmed slot booking.
    \item \textbf{Availability data accuracy:} agreement between displayed status and actual station status at arrival.
    \item \textbf{Station utilization:} change in booked-slot usage for participating stations.
    \item \textbf{User engagement:} repeat-usage rate and EcoCoins accumulation per active user.
\end{itemize}

\subsection{Pilot validation plan %
\BvrHint{\ldots high-level}}
\begin{itemize}[leftmargin=0.7cm]
    \item Recruit 10--20 pilot EV users and coordinate with 2--3 charging station operators (simulated if real operators are unavailable).
    \item Compare metrics before/after within the iteration window by logging baseline estimates via short interviews or existing walk-in behaviour.
\end{itemize}

\section{Scalability %
\BvrHint{\ldots with foresight}}
The proposition should scale in both theoretical and
practical senses.

\begin{enumerate}[leftmargin=0.7cm]
    \item \textbf{Scalable (theoretically):} the system should support growth in the number of stations, bookings, and concurrent users by:
    \begin{itemize}[leftmargin=0.7cm]
        \item decoupling components (frontend/backend),
        \item enabling pagination and indexing for common location queries,
        \item using a stateless, event-driven API design.
    \end{itemize}
    
    \item \textbf{Operationally deployable (practically):} the system should run on common web hosting:
    \begin{itemize}[leftmargin=0.7cm]
        \item cloud-native database (MongoDB Atlas),
        \item containerized or platform-hosted deployment,
        \item CI/CD pipeline for staging/production.
    \end{itemize}
\end{enumerate}

\section{Engine availability heuristic}
A mature set of ``engines'' (tooling/services) is available
to implement this as a web-based project:
\begin{itemize}[leftmargin=0.7cm]
    \item Web framework for REST APIs (Node.js, Express.js).
    \item Managed cloud document database (MongoDB Atlas).
    \item Location-intelligence APIs (Google Maps, Places, Directions) for search, geocoding, and routing.
    \item Authentication approach (token-based or session-based).
    \item Test framework for unit/integration tests.
    \item CI runner and staging deployment target.
\end{itemize}
We will use these mature components to focus on workflow
design, real-time correctness, and measurement rather than
inventing new infrastructure.

\section{Project scope and deliverables %
\BvrHint{\ldots iteration-friendly}}
\subsection{Initial deliverable %
\BvrHint{\ldots first iteration}}
\begin{itemize}[leftmargin=0.7cm]
    \item Centralized EV charging station locator (search + map view)
    \item Real-time availability and status dashboard
    \item Online slot booking \& reservation flow
    \item Smart search \& filters (connector type, distance, availability, price)
    \item CI: build + backend unit tests + linting
    \item CD to staging with a single command or pipeline job
\end{itemize}

\subsection{Subsequent deliverables}
\begin{itemize}[leftmargin=0.7cm]
    \item Trip planning with charging stops for long-distance routes
    \item Rewards, badges \& eco-friendly impact tracking (EcoCoins)
    \item Admin dashboard for station operators (utilization, schedules)
    \item Improved search/filtering and indexed queries for scale readiness
\end{itemize}

\section{Risks and mitigations}
\begin{itemize}[leftmargin=0.7cm]
    \item \textbf{Stale or incorrect availability data:} rely on periodic sync with station status and clearly timestamp last-updated status; flag stale data in the UI.
    \item \textbf{Third-party API costs/limits:} cache Google Maps/Places/Directions responses where possible and monitor quota usage.
    \item \textbf{Low initial station coverage:} start with a curated set of stations/operators for the pilot and expand incrementally.
    \item \textbf{User adoption:} keep the booking flow to a minimal number of steps and use the EcoCoins/rewards loop to encourage repeat use.
    \item \textbf{Location data privacy:} minimize stored location data; only persist what is needed for search and trip planning.
\end{itemize}

\section{Summary}
This project proposes ChargEV, a deployable web platform
that eliminates range anxiety by making EV charging
simple, convenient, and predictable through a centralized
locator, real-time availability, and smart online booking.
It meets the course criteria by emphasizing time-to-value
through rapid iterations, implementing CI/CD to support
frequent safe releases, and using measurable evaluation
criteria (especially time-to-charge and booking conversion).
It remains focused on engineering workflow, real-time data
integration, and scalability readiness rather than
research-driven novelty.

\end{document}
